% Solution: Gemini / Official Hybrid (Problem Sheet 20, Frobenius Inner Product)
\begin{exercisesolution}[Solution to \cref{exc:frobenius_hilbert_schmidt_inner_product}]
\label{sol:frobenius_hilbert_schmidt_inner_product}
Let $V = \M_{n \times n}(\mathbb{R})$ and $\langle A, B \rangle := \Tr(A\transp B)$.
\begin{enumerate}[label=\textbf{(\alph*)}]
    \item \textbf{Inner product axioms:}
    \begin{itemize}
        \item \emph{Bilinearity:} For $A_1, A_2, B \in V$ and $\alpha \in \mathbb{R}$:
        \[
        \langle A_1 + \alpha A_2, B \rangle = \Tr\big((A_1 + \alpha A_2)\transp B\big) = \Tr(A_1\transp B) + \alpha \Tr(A_2\transp B) = \langle A_1, B \rangle + \alpha \langle A_2, B \rangle,
        \]
        using linearity of transpose and trace. Linearity in the second argument follows similarly.
        \item \emph{Symmetry:} Since $\Tr(M\transp) = \Tr(M)$ for any square matrix $M$:
        \[
        \langle A, B \rangle = \Tr(A\transp B) = \Tr\big((A\transp B)\transp\big) = \Tr(B\transp A) = \langle B, A \rangle.
        \]
        \item \emph{Positivity:} Computing the trace of $A\transp A$:
        \[
        \langle A, A \rangle = \Tr(A\transp A) = \sum_{i=1}^n (A\transp A)_{ii} = \sum_{i=1}^n \sum_{j=1}^n (A\transp)_{ij} A_{ji} = \sum_{i=1}^n \sum_{j=1}^n a_{ji}^2 = \sum_{i,j=1}^n a_{ij}^2 \ge 0.
        \]
        Moreover, $\langle A, A \rangle = \sum_{i,j=1}^n a_{ij}^2 = 0$ if and only if $a_{ij} = 0$ for all $i,j$, i.e., $A = 0$.
    \end{itemize}

    \item \textbf{Orthonormal basis:}
    The standard elementary matrices $\{E_{ij} \mid 1 \le i, j \le n\}$, where $(E_{ij})_{k\ell} = \delta_{ik}\delta_{j\ell}$, form an orthonormal basis because:
    \[
    \langle E_{ij}, E_{k\ell} \rangle = \Tr(E_{ij}\transp E_{k\ell}) = \Tr(E_{ji} E_{k\ell}) = \Tr(\delta_{ik} E_{j\ell}) = \delta_{ik} \Tr(E_{j\ell}) = \delta_{ik} \delta_{j\ell}.
    \]

    \item \textbf{Hilbert-Schmidt norm formula:}
    The induced norm of $A \in V$ is:
    \[
    \|A\|_{HS} = \sqrt{\langle A, A \rangle} = \sqrt{\Tr(A\transp A)} = \sqrt{\sum_{i=1}^n \sum_{j=1}^n a_{ij}^2}. \qedhere
    \]
\end{enumerate}
\exinfo{Hybrid Solution (Gemini 3.7 Flash / Biran Lecture Notes). Verifies the Frobenius / Hilbert-Schmidt inner product axioms and orthonormal basis on $\M_{n \times n}(\mathbb{R})$.}
\end{exercisesolution}

% Solution: Gemini / Official Hybrid (Problem Sheet 20, Adjoint Matrix Properties)
\begin{exercisesolution}[Solution to \cref{exc:properties_adjoint_matrix_rules}]
\label{sol:properties_adjoint_matrix_rules}
Let $A, B \in \M_{n \times m}(K)$ and $C \in \M_{m \times p}(K)$.
\begin{enumerate}[label=\textbf{(\alph*)}]
    \item \textbf{Additivity:}
    $\big((A + B)^*\big)_{ij} = \overline{(A + B)_{ji}} = \overline{A_{ji} + B_{ji}} = \overline{A_{ji}} + \overline{B_{ji}} = (A^*)_{ij} + (B^*)_{ij} = (A^* + B^*)_{ij}$.

    \item \textbf{Conjugate-homogeneity:}
    $\big((\lambda A)^*\big)_{ij} = \overline{(\lambda A)_{ji}} = \overline{\lambda A_{ji}} = \overline{\lambda} \, \overline{A_{ji}} = \overline{\lambda} (A^*)_{ij} = (\overline{\lambda} A^*)_{ij}$.

    \item \textbf{Involution:}
    $\big((A^*)^*\big)_{ij} = \overline{(A^*)_{ji}} = \overline{\overline{A_{ij}}} = A_{ij}$, so $(A^*)^* = A$.

    \item \textbf{Identity:}
    $(I_n^*)_{ij} = \overline{(I_n)_{ji}} = \overline{\delta_{ji}} = \delta_{ij} = (I_n)_{ij}$.

    \item \textbf{Product rule:}
    \begin{align*}
    \big((A C)^*\big)_{ij} &= \overline{(A C)_{ji}} = \overline{\sum_{k=1}^m A_{jk} C_{ki}} = \sum_{k=1}^m \overline{A_{jk}} \, \overline{C_{ki}} = \sum_{k=1}^m \overline{C_{ki}} \, \overline{A_{jk}} \\
    &= \sum_{k=1}^m (C^*)_{ik} (A^*)_{kj} = (C^* A^*)_{ij}. \qedhere
    \end{align*}
\end{enumerate}
\exinfo{Hybrid Solution (Gemini 3.7 Flash / Biran Lecture Notes). Establishes fundamental algebraic properties of the conjugate transpose adjoint on $\M_{n \times m}(\mathbb{C})$.}
\end{exercisesolution}

% Solution: Gemini / Official Hybrid (Problem Sheet 20, Matrix Types and Subspaces)
\begin{exercisesolution}[Solution to \cref{exc:hermitian_trace_subspaces_sc}]
\label{sol:hermitian_trace_subspaces_sc}
\begin{enumerate}[label=\textbf{(\alph*)}]
    \item Let $A = \begin{pmatrix} x & -x \\ x & x \end{pmatrix}$. Then:
    \[
    A^* A = \begin{pmatrix} \overline{x} & \overline{x} \\ -\overline{x} & \overline{x} \end{pmatrix} \begin{pmatrix} x & -x \\ x & x \end{pmatrix} = \begin{pmatrix} |x|^2 + |x|^2 & -\overline{x}x + \overline{x}x \\ -\overline{x}x + \overline{x}x & |-\overline{x}|^2 + |x|^2 \end{pmatrix} = \begin{pmatrix} 2|x|^2 & 0 \\ 0 & 2|x|^2 \end{pmatrix}.
    \]
    Thus $A^* A = I_2 \iff 2|x|^2 = 1 \iff |x| = \frac{1}{\sqrt{2}}$.
    The matrix $A$ is unitary for all $x \in \mathbb{C}$ with $|x| = \frac{1}{\sqrt{2}}$.

    \item \textbf{$\mathbb{C}$-linear subspace test:}
    \begin{itemize}
        \item \emph{Symmetric matrices:} Is a $\mathbb{C}$-linear subspace: if $A\transp = A$ and $B\transp = B$, then $(\alpha A + \beta B)\transp = \alpha A\transp + \beta B\transp = \alpha A + \beta B$ for all $\alpha, \beta \in \mathbb{C}$.
        \item \emph{Hermitian matrices:} Not a $\mathbb{C}$-linear subspace: if $A = I_n$ (Hermitian), then $(i I_n)^* = -i I_n^* = -i I_n \ne i I_n$, so it is not closed under complex scalar multiplication.
        \item \emph{Unitary matrices:} Not a subspace: the zero matrix is not unitary ($0^* 0 = 0 \ne I_n$).
        \item \emph{Normal matrices:} Not a subspace: the sum of two normal matrices need not be normal (e.g. $A = \begin{pmatrix} 1 & 0 \\ 0 & 0 \end{pmatrix}$ and $B = \begin{pmatrix} 0 & 1 \\ 0 & 0 \end{pmatrix} + \begin{pmatrix} 0 & 0 \\ 1 & 0 \end{pmatrix}$).
    \end{itemize}

    \item \textbf{Real trace and determinant of Hermitian matrices:}
    Let $A^* = A$.
    \begin{itemize}
        \item Trace: $\overline{\Tr(A)} = \Tr(\overline{A}) = \Tr(\overline{A}\transp) = \Tr(A^*) = \Tr(A)$, which implies $\Tr(A) \in \mathbb{R}$.
        \item Determinant: $\overline{\det(A)} = \det(\overline{A}) = \det(\overline{A}\transp) = \det(A^*) = \det(A)$, which implies $\det(A) \in \mathbb{R}$. \qedhere
    \end{itemize}
\end{enumerate}
\exinfo{Hybrid Solution (Gemini 3.7 Flash / Biran Lecture Notes). Analyzes unitarity conditions, subspace closure for matrix types, and reality of Hermitian traces/determinants.}
\end{exercisesolution}

% Solution: Gemini / Official Hybrid (Problem Sheet 20, Gram Matrix)
\begin{exercisesolution}[Solution to \cref{exc:gram_matrix_properties}]
\label{sol:gram_matrix_properties}
Let $A \in \M_{n \times n}(\mathbb{R})$.
\begin{enumerate}[label=\textbf{(\alph*)}]
    \item $(A\transp A)\transp = A\transp (A\transp)\transp = A\transp A$, so $A\transp A$ is symmetric.

    \item For any $v \in \mathbb{R}^n$, $v\transp (A\transp A) v = (Av)\transp (Av) = \|Av\|^2 \ge 0$.
    Thus $v\transp (A\transp A) v > 0$ for all $v \ne 0 \iff \|Av\| > 0$ for all $v \ne 0 \iff Av \ne 0$ for all $v \ne 0 \iff \ker A = \{0\} \iff A$ is invertible.

    \item We show $\ker(A\transp A) = \ker(A)$:
    \begin{itemize}
        \item If $v \in \ker(A)$, then $Av = 0 \implies A\transp A v = A\transp 0 = 0$, so $v \in \ker(A\transp A)$.
        \item Conversely, if $v \in \ker(A\transp A)$, then $A\transp A v = 0 \implies \|Av\|^2 = v\transp A\transp A v = v\transp 0 = 0 \implies Av = 0$, so $v \in \ker(A)$.
    \end{itemize}
    By the Rank-Nullity Theorem:
    \[
    \operatorname{rank}(A\transp A) = n - \dim\ker(A\transp A) = n - \dim\ker(A) = \operatorname{rank}(A). \qedhere
    \]
\end{enumerate}
\exinfo{Hybrid Solution (Gemini 3.7 Flash / Biran Lecture Notes). Proves symmetry, positive definiteness, and rank preservation of the Gram matrix $A\transp A$.}
\end{exercisesolution}

% Solution: Gemini / Official Hybrid (Problem Sheet 20, Laguerre Weight Inner Product)
\begin{exercisesolution}[Solution to \cref{exc:exponential_weighted_polynomial_inner_product}]
\label{sol:exponential_weighted_polynomial_inner_product}
\begin{enumerate}[label=\textbf{(\alph*)}]
    \item \textbf{Scalar product verification:}
    \begin{itemize}
        \item \emph{Bilinearity and Symmetry:} Follow immediately from the linearity of improper Riemann integrals and commutativity of polynomial multiplication: $p(t)q(t)e^{-t} = q(t)p(t)e^{-t}$.
        \item \emph{Positivity:} For any non-zero polynomial $p \in \mathbb{R}_{\le n}[x]$, the integrand $p(t)^2 e^{-t}$ is continuous and non-negative on $[0, \infty)$.
        Since $p$ has at most $n$ real roots, there exists an interval $[a, b] \subset (0, \infty)$ with $a < b$ on which $p(t)^2 > 0$.
        Therefore:
        \[
        \langle p, p \rangle = \int_0^\infty p(t)^2 e^{-t} \, dt \ge \int_a^b p(t)^2 e^{-t} \, dt > 0.
        \]
    \end{itemize}

    \item \textbf{Matrix representation:}
    In the standard basis $\mathcal{B} = (1, x, \dots, x^n)$, the $(i, j)$-entry for $0 \le i, j \le n$ is given by the Gamma function evaluation:
    \[
    a_{ij} = \langle x^i, x^j \rangle = \int_0^\infty t^{i+j} e^{-t} \, dt = (i + j)!.
    \]
    Thus the representation matrix is the Hankel matrix:
    \[
    [\langle \cdot, \cdot \rangle]_\mathcal{B} = \begin{pmatrix}
      0! & 1! & 2! & \dots & n! \\
      1! & 2! & 3! & \dots & (n+1)! \\
      2! & 3! & 4! & \dots & (n+2)! \\
      \vdots & \vdots & \vdots & \ddots & \vdots \\
      n! & (n+1)! & (n+2)! & \dots & (2n)!
    \end{pmatrix}. \qedhere
    \]
\end{enumerate}
\exinfo{Hybrid Solution (Gemini 3.7 Flash / Biran Lecture Notes). Verifies exponential-weighted polynomial inner product and computes its factorial Hankel Gram matrix.}
\end{exercisesolution}

% Solution: Gemini / Official Hybrid (Problem Sheet 20, Jordan-von Neumann Theorem)
\begin{exercisesolution}[Solution to \cref{exc:jordan_von_neumann_parallelogram}]
\label{sol:jordan_von_neumann_parallelogram}
\begin{enumerate}[label=\textbf{(\alph*)}]
    \item \textbf{Jordan--von Neumann Theorem:}
    \begin{itemize}
        \item \emph{($\Rightarrow$):} If $\|\cdot\|$ is induced by an inner product $\langle \cdot, \cdot \rangle$, then:
        \[
        \|x+y\|^2 + \|x-y\|^2 = \big(\|x\|^2 + 2\langle x, y \rangle + \|y\|^2\big) + \big(\|x\|^2 - 2\langle x, y \rangle + \|y\|^2\big) = 2\|x\|^2 + 2\|y\|^2.
        \]
        \item \emph{($\Leftarrow$):} Define $\langle x, y \rangle := \frac{1}{4}\big(\|x+y\|^2 - \|x-y\|^2\big)$.
        Symmetry is clear: $\langle y, x \rangle = \frac{1}{4}(\|y+x\|^2 - \|y-x\|^2) = \langle x, y \rangle$.
        Positivity holds since $\langle x, x \rangle = \frac{1}{4}\|2x\|^2 = \|x\|^2 \ge 0$, with equality only for $x = 0$.
        Additivity in the first variable: applying the parallelogram identity to $(x+z, y)$ and $(x-z, y)$ shows:
        \[
        \langle x+z, y \rangle + \langle x-z, y \rangle = 2\langle x, y \rangle.
        \]
        Setting $z = x$ gives $\langle 2x, y \rangle = 2\langle x, y \rangle$, and setting $x = z = \frac{u+v}{2}$ yields additivity $\langle u+v, y \rangle = \langle u, y \rangle + \langle v, y \rangle$.
        By induction $\langle nx, y \rangle = n\langle x, y \rangle$ for $n \in \mathbb{N}$, extending to $\mathbb{Q}$ by division, and to $\mathbb{R}$ by continuity of the norm.
    \end{itemize}

    \item \textbf{The $\ell_1$-norm:}
    $\|v\|_1 = \sum_{i=1}^n |v_i|$ satisfies:
    \begin{itemize}
        \item $\|v\|_1 \ge 0$, and $\|v\|_1 = 0 \iff |v_i| = 0 \ \forall i \iff v = 0$.
        \item $\|\alpha v\|_1 = \sum |\alpha v_i| = |\alpha| \sum |v_i| = |\alpha| \|v\|_1$.
        \item $\|u + v\|_1 = \sum |u_i + v_i| \le \sum (|u_i| + |v_i|) = \|u\|_1 + \|v\|_1$.
    \end{itemize}
    Thus $\|\cdot\|_1$ is a valid norm.
    For $n \ge 2$, take $x = e_1 = (1, 0, \dots, 0)$ and $y = e_2 = (0, 1, \dots, 0)$.
    Then $\|x\|_1 = 1, \|y\|_1 = 1, \|x+y\|_1 = 2, \|x-y\|_1 = 2$.
    The parallelogram identity evaluates to:
    \[
    \|x+y\|_1^2 + \|x-y\|_1^2 = 2^2 + 2^2 = 8 \ne 4 = 2(1^2) + 2(1^2) = 2\|x\|_1^2 + 2\|y\|_1^2.
    \]
    Thus $\|\cdot\|_1$ is not induced by any scalar product for $n \ge 2$. \qedhere
\end{enumerate}
\exinfo{Hybrid Solution (Gemini 3.7 Flash / Biran Lecture Notes). Proves the Jordan-von Neumann parallelogram characterization of inner product spaces and tests the $\ell_1$-norm.}
\end{exercisesolution}

% Solution: Gemini / Official Hybrid (Problem Sheet 20, Rotation Invariance)
\begin{exercisesolution}[Solution to \cref{exc:rotation_invariance_angle}]
\label{sol:rotation_invariance_angle}
Let $V = \mathbb{R}^2$ with the standard scalar product $\langle u, v \rangle = u\transp v$.
A planar rotation by an angle $\theta$ is represented in the standard basis by the orthogonal matrix:
\[
R_\theta = \begin{pmatrix} \cos\theta & -\sin\theta \\ \sin\theta & \cos\theta \end{pmatrix}.
\]
Direct calculation shows:
\[
R_\theta\transp R_\theta = \begin{pmatrix} \cos\theta & \sin\theta \\ -\sin\theta & \cos\theta \end{pmatrix} \begin{pmatrix} \cos\theta & -\sin\theta \\ \sin\theta & \cos\theta \end{pmatrix} = \begin{pmatrix} \cos^2\theta + \sin^2\theta & 0 \\ 0 & \sin^2\theta + \cos^2\theta \end{pmatrix} = I_2.
\]
Therefore, for any $v_1, v_2 \in \mathbb{R}^2$:
\[
\langle R_\theta v_1, R_\theta v_2 \rangle = (R_\theta v_1)\transp (R_\theta v_2) = v_1\transp (R_\theta\transp R_\theta) v_2 = v_1\transp I_2 v_2 = v_1\transp v_2 = \langle v_1, v_2 \rangle.
\]
In particular, $\|R_\theta v\| = \sqrt{\langle R_\theta v, R_\theta v \rangle} = \sqrt{\langle v, v \rangle} = \|v\|$.
Substituting these into the angle formula yields:
\[
\cos \angle(R_\theta v_1, R_\theta v_2) = \frac{\langle R_\theta v_1, R_\theta v_2 \rangle}{\|R_\theta v_1\| \|R_\theta v_2\|} = \frac{\langle v_1, v_2 \rangle}{\|v_1\| \|v_2\|} = \cos \angle(v_1, v_2). \qedhere
\]
\exinfo{Hybrid Solution (Gemini 3.7 Flash / Biran Lecture Notes). Demonstrates preservation of scalar products, norms, and angles under orthogonal planar rotations.}
\end{exercisesolution}

% Official Solution (Problem Sheet 23, Exercise 1)
\begin{exercisesolution}[Solution to \cref{exc:extend_orthonormal_system_sheet23}]
\label{sol:extend_orthonormal_system_sheet23}
Let $S = (u_1, \dots, u_k)$ be an orthonormal system in a finite-dimensional Euclidean space $V$ with $n := \dim V$.
By \cref{thm:properties_orthogonal_systems} \textbf{(a)}, $S$ is linearly independent.
By the Steinitz exchange lemma (\cref{lem:strong_steinitz}), we can extend $S$ to a basis $\mathcal{B} = (u_1, \dots, u_k, v_{k+1}, \dots, v_n)$ of $V$.
We apply the Gram-Schmidt orthogonalization process (\cref{thm:gram_schmidt}) to $\mathcal{B}$.
For $j \le k$, since $(u_1, \dots, u_k)$ is already orthonormal, the projection formulas give:
\[
w_j = u_j - \sum_{i=1}^{j-1} \langle u_j, u_i \rangle u_i = u_j - 0 = u_j, \qquad \|w_j\| = \|u_j\| = 1.
\]
Thus, the first $k$ vectors are completely unchanged: $e_j = u_j$ for $1 \le j \le k$.
The remaining vectors $e_{k+1}, \dots, e_n$ produced by Gram-Schmidt are orthogonal to each other and to $u_1, \dots, u_k$, and each has unit norm.
Thus, $(u_1, \dots, u_k, e_{k+1}, \dots, e_n)$ is an orthonormal basis of $V$ extending $S$.
\exinfo{Official Solution (Problem Sheet 23, Exercise 1). Proves that any orthonormal system in a finite-dimensional Euclidean space extends to an orthonormal basis via Gram-Schmidt.}
\end{exercisesolution}

% Official Solution (Problem Sheet 20, Exercise 4)
\begin{exercisesolution}[Solution to \cref{exc:gram_schmidt_polynomials_derivative}]
\label{sol:gram_schmidt_polynomials_derivative}
\begin{enumerate}[label=\textbf{(\alph*)}]
    \item \textbf{Gram-Schmidt process on $(1, x, x^2)$:}
    Let $b_1 = 1, b_2 = x, b_3 = x^2$.
    \begin{itemize}
        \item First vector: $w_1 = 1$, $\|w_1\|^2 = \int_0^1 1\,dx = 1 \implies u_1 = 1$.
        \item Second vector: $\langle b_2, u_1 \rangle = \int_0^1 x\,dx = \frac{1}{2}$.
        \[
        w_2 = x - \frac{1}{2} u_1 = x - \frac{1}{2}.
        \]
        $\|w_2\|^2 = \int_0^1 (x - \frac{1}{2})^2 \, dx = \left[ \frac{(x-1/2)^3}{3} \right]_0^1 = \frac{1/8 - (-1/8)}{3} = \frac{1}{12}$.
        Thus $u_2 = \frac{w_2}{\|w_2\|} = \sqrt{12}(x - \frac{1}{2}) = \sqrt{3}(2x - 1)$.
        \item Third vector:
        $\langle b_3, u_1 \rangle = \int_0^1 x^2\,dx = \frac{1}{3}$.
        $\langle b_3, u_2 \rangle = \sqrt{3} \int_0^1 x^2(2x - 1)\,dx = \sqrt{3} \left[ \frac{2x^4}{4} - \frac{x^3}{3} \right]_0^1 = \sqrt{3}(\frac{1}{2} - \frac{1}{3}) = \frac{\sqrt{3}}{6}$.
        \[
        w_3 = x^2 - \frac{1}{3}(1) - \frac{\sqrt{3}}{6} \cdot \sqrt{3}(2x - 1) = x^2 - \frac{1}{3} - \frac{1}{2}(2x - 1) = x^2 - x + \frac{1}{6}.
        \]
        $\|w_3\|^2 = \int_0^1 (x^2 - x + \frac{1}{6})^2 \, dx = \int_0^1 (x^4 - 2x^3 + \frac{4}{3}x^2 - \frac{1}{3}x + \frac{1}{36})\,dx = \frac{1}{5} - \frac{2}{4} + \frac{4}{9} - \frac{1}{6} + \frac{1}{36} = \frac{1}{180}$.
        Thus $u_3 = \sqrt{180}(x^2 - x + \frac{1}{6}) = \sqrt{5}(6x^2 - 6x + 1)$.
    \end{itemize}
    The orthonormal basis is $\mathcal{C} = \big(1, \ \sqrt{3}(2x-1), \ \sqrt{5}(6x^2 - 6x + 1)\big)$.

    \item \textbf{Upper triangular representation matrix for derivative:}
    The basis $\mathcal{C} = (u_1, u_2, u_3)$ satisfies $\deg(u_k) = k - 1$.
    Because differentiation lowers polynomial degree:
    \begin{align*}
    D(u_1) &= 0, \\
    D(u_2) &= 2\sqrt{3} = 2\sqrt{3} u_1 \in \Sp(u_1), \\
    D(u_3) &= \sqrt{5}(12x - 6) = 6\sqrt{5}(2x - 1) = 2\sqrt{15} u_2 \in \Sp(u_1, u_2).
    \end{align*}
    Hence the representation matrix $[D]_\mathcal{C}^\mathcal{C} = \begin{pmatrix} 0 & 2\sqrt{3} & 0 \\ 0 & 0 & 2\sqrt{15} \\ 0 & 0 & 0 \end{pmatrix}$ is strictly upper triangular. \qedhere
\end{enumerate}
\exinfo{Official Solution (Problem Sheet 20, Exercise 4). Computes orthonormal Legendre-type polynomials on $[0,1]$ via Gram-Schmidt and derives upper triangular derivative matrix.}
\end{exercisesolution}
